\scanpage{5}
Donc pour que ces trois points forment un triangle équilatéral (il est de toutes façons \unsure{isoscèle|isocèle}) il fst que l'on ait
\[
\deleted{\unsure{4}}\cos\frac{\alpha}{2}=\pm\frac{1}{2\sin\theta/2}\quad\deleted{\unsure{\sin}}\ .\tag*{\textcircled{1}}
\]
\margin{[\uline{NB} Un tel $\alpha$ existe sss $|\sin\frac{\theta}{2}|\geq\frac12$, i.e. si $0\leq\theta\leq\pi$, sss \inserted{sss} $\frac{\pi}{6}\leq\frac{\theta}{2}\leq\frac{5\pi}{6}$ i.e. $\deleted{\unsure{\pi/3|\frac{\pi}{3}\leq}}\boxed{\theta\geq\frac{\pi}{3}}$]}
ou encore, comme $\cos\alpha=\cos^2\frac{\alpha}{2}-\sin^2\frac{\alpha}{2}=2\cos^2\frac{\alpha}{2}-1$
\[
\cos\alpha=\frac{1}{2\sin^2\theta/2}-1=\deleted{\unsure{\cos\theta}}\frac{1-2\sin^2\theta/2}{2\sin^2\theta/2}
=\frac{\cos\deleted{\unsure{2}}\theta}{2\sin^2\theta/2}=\frac{\cos\theta}{1-\cos\theta}\tag*{\textcircled{$1'$}}
\]
% slanted
\margin{\uline{NB} $\left\{\begin{aligned}\cos2\alpha&=-\frac35\\\cos\alpha&=\frac{1}{\sqrt5}\end{aligned}\right.$ pour $\theta=\frac{2\pi}{5}$}

Bien entendu, dans ce cas, les autres ``angles dièdres'' définis par ce triangle en les deux autres sommets, sont également égaux à $\theta$. On a en fait construit ici, \inserted{choisissant la détermination principale de $\alpha$ comme $0<\alpha<\pi$,} \uline{le} seul triangle équilatéral inscrit dans la sphère $S$, d'angle dièdre $\theta$, ayant un sommet en $u$, un \deleted{\unsure{côté|autre}} \inserted{deuxième sommet $x$} dans la direction tangente $v$ à $S$ en $u$, et \deleted{situé} tel qu'avec le troisième sommet $x'$ l'orientation $u,x,x'$ soit directe.

Nous allons fixer par la suite
\[
\theta=\frac{2\pi}{5}\tag*{\textcircled{2}}
\]
et, pour quelque temps, on fixe $u$ et $v$. On posera
\[
v_i=v\cos i\theta+w\sin i\theta\qquad i\in\mathbb Z/5\mathbb Z\tag*{(3)}
\]
(de sorte que $v_0=v$, i.e. $v_i=v$ si $i\equiv0\ (5)$), et on posera
\[
\deleted{\unsure{v_i=v\ \text{si }i\equiv0\ (5)}}
\]
\[
\begin{aligned}
x_i&=x(u,v_i,\alpha)=u\cos\alpha+v_i\sin\alpha
% Both occurrences of sin alpha carry a transposition stroke to before the preceding factor.
=u\cos\alpha+v\sin\alpha\cos i\theta+w\sin\alpha\sin i\theta\\
&& (i\in\mathbb Z/5\mathbb Z)
\end{aligned}\tag*{(4)}
\]
\margin{\uline{NB} Il est immédiat que les $u,x_i$ sont $6$ pts \uline{distincts} (leurs projections sur le plan $u^\perp$ sont $0$ et les $v_i\sin\alpha$) à coordonnée suivant $u>0$ ; donc les $u,u'$, $x_i',x_i$ sont $12$ points \uline{distincts} de $S$}

où $\alpha$ est fixé par (1) ou (1$'$), $0<\alpha<\pi$. On trouve donc cinq points en ordre circulaire sur le cercle de \inserted{centre $u$} rayon (rectiligne) $4\sin^2\frac{\alpha}{2}$ tracé sur la sphère $S$. On pose de plus
\[
x'_i=-x_i,\quad u'=-u\qquad (i\in\mathbb Z/5\mathbb Z)\tag*{(5)}
\]
\margin{On appelle cet ens de 12 pts l'\uline{icosaèdre} défini \uline{par} $u,v$.}

Un ensemble de points $(x,y,z)$ de $S$ \inserted{définissant un triangle équilatéral} d'angle au centre $\alpha$ est appelé un \uline{triangle icosaèdral} ; \inserted{cela signifie aussi qu'il y a une rotation qui l'amène en $(u,x_0,x_1)$} un ensemble de points $(a,b_0,b_1,\ldots,b_4)$ \inserted{de $S$} est appelé un \uline{calotte d'icosaèdre} s'il existe une transformation \inserted{orthogonale} $g\in O(E)$ qui l'amène en l'ensemble précédent $(u,x_0,x_1,\ldots,x_4)$ ; Il revient encore au même \deleted{que} de dire \inserted{son cardinal est $6$, et que l'on peut ordonner cet ens. de telle façon que} que les triangles $(a,b_i,b_{i+1})$ pour $0\leq i\leq3$ sont icosaèdraux (et alors $(a,b_4,b_0)$ est aussi icosaèdral$^{*}$. Le point $a$ est alors uniquement
% archive sheet number: 4
